% ====================================================================
% Section 10: Abstract and Conclusion
% ====================================================================
% This file intentionally contains both the manuscript abstract and the
% concluding section so the final paper assembly can move the abstract to
% the front matter while keeping the argument synchronized with the close.

\begin{abstract}
Agent self-evolution reframes large language model agents not as static prompt-response systems, but as adaptive software organisms whose prompts, memories, tools, policies, workflows, evaluators, and code can be repeatedly varied, tested, selected, and retained. This survey synthesized evidence from 107 papers, 348 repository records, 97 user pain-point reports, 1306 blog and technical sources, a 19+ project architecture-search index, and a researcher network spanning academic, industrial, and open-source communities. Across these materials, we argued that the field is converging on five recurring evolutionary loops: verifier- and reward-driven improvement, self-play and self-generated curriculum, reflection and context evolution, architecture and code self-modification, and memory- or experience-centered lifelong adaptation. Hybrid systems combine these loops, but the loops themselves explain why work as different as STaR, Reflexion, Voyager, Self-Rewarding LM, ADAS, DGM, AlphaEvolve, ACE, ReasoningBank, and AriadneMem can be studied under one framework.

The central conclusion is deliberately double-sided. On the positive side, agent evolution is already producing a shared cross-domain vocabulary for machine learning, evolutionary computation, AutoML, software engineering, cognitive architectures, and human-AI organization: generate candidates, expose them to feedback, select under constraints, preserve lineage, and reuse experience. On the cautionary side, the same mechanism makes evaluation, governance, and safety central rather than secondary. Reward hacking, benchmark overfitting, capability erosion, experience poisoning, runaway cost, and unsafe self-modification are not deployment details; they are failure modes of evolution itself. The practical path forward is therefore not unconstrained recursive self-improvement, but governed open-endedness: verifiable tasks, independent evaluators, sandboxed mutation, auditable memory, human escalation for high-risk changes, and cross-domain tests that measure transfer rather than leaderboard movement alone. If those infrastructures mature, agent self-evolution can become an engineering discipline for building systems that learn from interaction while remaining inspectable, reversible, and socially accountable.
\end{abstract}

\section{Conclusion: Toward Governed Open-Ended Agent Evolution}
\label{sec:conclusion}

\subsection{Returning to the Question Posed in the Introduction}

The introduction opened with a shift in the basic question of agent research. The crucial question is no longer only how to make a model answer better in a single invocation. It is how to make an agentic system improve through time: after it acts, fails, receives feedback, stores experience, modifies some part of itself, and re-enters the world with a verifiably changed capability profile. That shift is the organizing claim of this survey, and the evidence reviewed across theory, methods, systems, evaluation, industry practice, user pain points, and future directions supports it. Agent self-evolution is not a single algorithm, a single benchmark, or a single architecture. It is a family of closed-loop designs in which mutable system components are exposed to selection pressure and then carried forward as reusable state.

This conclusion therefore returns to the four definitional requirements introduced at the beginning. First, an evolving agent must contain a mutable object: model weights, natural-language prompts, memory entries, tool-use policies, planning code, critic instructions, agent workflows, multi-agent topologies, data generators, or evaluation rules. Second, it must receive feedback: program tests, reward models, human preferences, peer debate, environment rewards, hidden benchmarks, cost measurements, safety checks, or long-horizon user outcomes. Third, it must have a selection or update mechanism: filtering, reinforcement learning, DPO, MCTS, text-gradient edits, archive selection, memory consolidation, code patching, workflow search, or curriculum generation. Fourth, its improvement must be auditable. A demo that appears smoother is not enough; the change must be traceable through scores, ablations, regression tests, lineage records, cost logs, failure analyses, or safety reviews.

The survey also confirmed the introduction's broader motivation. Static prompting is mismatched with tasks that require exploration, failure recovery, and adaptation. Human supervision and hand-built workflows are too expensive to cover long-tail environments. Real applications require memory across sessions, adaptation to changing tools and codebases, and robustness under distribution shift. At the same time, self-modifying systems create new safety obligations because the object being optimized may include the very instructions, memory, and code that determine future behavior. The result is a field pulled by two forces: the promise of compounding improvement, and the need to govern that compounding before it becomes brittle, deceptive, expensive, or unsafe.

\subsection{The Five Evolutionary Loops}

Although the literature contains many names and variants, this survey reduces the core mechanism to five evolutionary loops. The first is the verifier or reward loop. In this loop, the agent proposes solutions or trajectories, a verifier evaluates them, and successful candidates are reinforced, stored, or used for further training. STaR demonstrated the early generate--filter--finetune pattern for reasoning; Self-Rewarding LM and Meta-Rewarding moved evaluation into model-generated judgments; RISE, Agent-R, RAGEN, and related RL systems transformed failures and trajectories into learning signals. The verifier loop is powerful because it grounds evolution in measurable outcomes, but it is also vulnerable to Goodhart's law: if the evaluator is incomplete, the agent can learn to satisfy the proxy rather than the task.

The second is the self-play and self-generated curriculum loop. Here the agent or a population of agents generates tasks, adversaries, games, or curricula that expose weaknesses and create new training pressure. Absolute Zero, SPIRAL, self-play fine-tuning, Multi-Agent Evolve, Agent0-style curricula, and other zero-data approaches all push toward reducing dependence on external datasets. This loop matters because open-ended intelligence cannot wait for humans to enumerate every challenge. Yet it also raises the question of task reality. A synthetic curriculum is useful only if it induces capabilities that transfer to real user distributions, real software systems, real scientific problems, or real operational constraints.

The third is the reflection, prompt, and context loop. Reflexion, Self-Refine, ExpeL, ACE, EvolveR, Agent Symbolic Learning, and related approaches show that many improvements can occur without changing model weights. Agents can critique their own failures, turn successful trajectories into rules, maintain playbooks, update prompts, or apply textual gradients to policies and roles. This loop is attractive because it is cheap, interpretable, and deployable on top of existing foundation models. Its weakness is that natural-language memory can become stale, contradictory, bloated, or vulnerable to injection. Context evolution must therefore be paired with compression, retrieval, provenance, and forgetting mechanisms.

The fourth is the architecture and code self-modification loop. ADAS treats agent design as a search problem. Gödel Agent, DGM, SICA, AlphaEvolve, FunSearch-style program evolution, EvoMAC, and architecture-search projects move the mutable object from prompt text to executable structure. This is the most ambitious loop because it asks the agent not only to choose actions, but to modify the machinery that chooses actions. It also has the highest governance burden. A code-changing agent needs sandboxing, version control, tests, rollback, permissions, and independent evaluators. Without those controls, code evolution becomes an uncontrolled software supply-chain event.

The fifth is the memory and experience loop. Voyager, Generative Agents, ReasoningBank, Memory-R1, AriadneMem, SEDM, ELL-style lifelong systems, and experience-bank projects show that long-term adaptation depends on turning trajectories into durable assets. Memory is not merely storage; it is a selection system over what the agent believes is worth carrying forward. The memory loop supports transfer, personalization, organizational learning, and cumulative skill formation. It also introduces new risks: privacy leakage, experience poisoning, capability erosion through bad consolidation, and irreversible accumulation of false beliefs. Mature self-evolving agents will need memory governance as carefully designed as model training governance.

These five loops also clarify why earlier six-way method lists often include a hybrid category. Hybrid systems are not a sixth primitive; they are compositions. WebEvolver combines world-model adaptation, task generation, and policy learning. DGM-style systems combine archive selection with code mutation. Production frameworks combine memory, workflow orchestration, human review, and evaluation dashboards. The field is moving from isolated loops toward loop orchestration: deciding which loop to activate, at what risk level, with what budget, under which evaluator, and with which rollback policy.

\subsection{What the 19+ Project and 348-Repository Analysis Adds}

The paper corpus explains the conceptual frontier, but the project corpus shows what the ecosystem is trying to build. The 19+ project architecture-search and self-evolution index, together with the wider 348-repository analysis, reveals a split between research-grade open-ended evolution and practitioner-grade self-improvement utilities. Projects such as ADAS, Gödel Agent, EvoAgent, EvoAgentX, DGM, GenericAgent, Agent0, Agent Zero, SE-Agent, AgentEvolver, STELLA, ELL-StuLife, WebRL, Hermes Self-Evolution, EvoMap Evolver, AutoResearchClaw, scientific code-evolution systems, Agent Symbolic Learning, MetaGPT, and DARWIN cover a broad design space: architecture search, runtime self-reference, skill crystallization, zero-data curricula, trajectory evolution, prompt and skill optimization, autonomous research, multi-agent organization, and adversarial evolution.

Several conclusions follow from this project-level view. First, Python remains the default language for experimental self-evolution because it gives direct access to LLM SDKs, tests, scripts, notebooks, and ML infrastructure. However, the presence of Node.js and product-oriented projects shows that evolution is entering user-facing toolchains and agent operating systems, not just research notebooks. Second, projects increasingly include docs, tests, Dockerfiles, CI, security files, and agent-operation guides. This suggests that the community is slowly moving from prompt demos toward software artifacts that can be inspected and maintained. Third, the most mature projects separate generation from evaluation. Whether the evaluator is a unit test, benchmark harness, LLM judge, task simulator, or human review step, the separation is what keeps improvement from collapsing into self-approval.

Fourth, the project corpus highlights a maturity ladder. At the lowest level are prompt optimizers and reflection wrappers: useful, inexpensive, but limited in structural change. Above them are workflow evolution systems that search over agent roles, tools, or coordination patterns. Above those are skill-evolution and memory systems that accumulate reusable assets. Above those are trajectory and RL-based systems that learn from many attempts. At the highest level are self-referential code and architecture systems that can rewrite parts of their own implementation. Each level increases capability, but also increases verification cost. A production team should not jump to code self-modification when a prompt, memory, or workflow loop is sufficient.

Finally, the repository analysis tempers hype. Star counts, trending status, and striking benchmark claims do not by themselves prove durable self-evolution. Many repositories are lists, tutorials, thin wrappers, or research prototypes. Others are impressive but expensive, hard to reproduce, or narrow to a benchmark. The project evidence therefore supports a pragmatic conclusion: the ecosystem has enough diversity to justify a unified survey, but not enough standardization to justify claims of a solved production discipline. The next stage is not merely more projects, but comparable project cards that report the mutable object, feedback source, evaluator independence, cost, safety boundary, lineage strategy, and transfer evidence.

\subsection{A Cross-Domain View: One Pattern, Many Substrates}

A central contribution of this survey is the cross-domain unification of agent evolution. In machine learning, the substrate is data, reward, model weights, and training objectives. In evolutionary computation, the substrate is a population of candidate programs or policies, an archive, a mutation operator, and a fitness function. In AutoML and NAS, the substrate is a pipeline, model architecture, or search space. In software engineering, the substrate is code, tests, version control, CI, and issue history. In cognitive architectures, the substrate is memory, goals, planning, reflection, and skill acquisition. In human organizations, the substrate is workflow, responsibility, review, escalation, and institutional memory.

Agent self-evolution sits at the intersection because modern LLM agents are hybrid artifacts. They are partly models, partly programs, partly documents, partly organizations, and partly interfaces to external tools. This is why single-discipline language is insufficient. A reinforcement-learning view explains reward optimization but not lineage governance. A software-engineering view explains tests and rollback but not open-ended exploration. An evolutionary-computation view explains diversity and stepping stones but not prompt injection or enterprise permissions. A cognitive-architecture view explains memory and reflection but not benchmark leakage. A human-organization view explains accountability but not automated search. The field needs all of them.

The cross-domain map also clarifies transfer. If an evolving coding agent learns only a benchmark-specific patching trick, it has not demonstrated general evolution. If a memory agent stores every failure but cannot retrieve the right lesson under a new task, it has not demonstrated lifelong learning. If an architecture-search system finds a complex multi-agent workflow that costs more than it saves, it has not demonstrated useful adaptation. Transfer requires abstract assets: validators, failure modes, tool-use templates, task decompositions, safety rules, memory schemas, code-edit patterns, and decision policies. These assets are the bridge from one domain to another.

This view changes how we should evaluate progress. The question is not simply whether a new system improves on SWE-Bench, WebArena, a math benchmark, a game score, or a customer-service task. The question is what kind of asset the improvement creates and whether that asset can be inspected, reused, composed, and safely updated. A benchmark gain that leaves no reusable knowledge may be less important than a smaller gain that produces a robust verifier, a reusable skill, or a clear failure taxonomy. Agent evolution should therefore be evaluated both by outcome metrics and by asset quality.

\subsection{The Researcher and Institution Map}

The researcher graph reviewed in this paper shows that agent self-evolution is not emerging from a single school. It is a convergence of several intellectual lineages. The STaR and Quiet-STaR line connects chain-of-thought reasoning, self-generated rationales, and self-training through researchers such as Eric Zelikman and Noah Goodman. The ReAct, Tree-of-Thoughts, and Reflexion line connects interactive reasoning, tool use, and verbal reinforcement through Shunyu Yao, Karthik Narasimhan, Noah Shinn, and collaborators. The Self-Refine and self-debugging line connects iterative critique with AI2, CMU, and Google code-generation work. The SPIN and self-play line connects game-theoretic alignment and self-generated preference pressure. The ADAS and DGM line connects Jeff Clune's open-ended evolution tradition with automated agent design and self-modifying code. The AlphaEvolve and FunSearch line connects DeepMind-style program search with LLM proposal generation and rigorous evaluation.

Institutionally, the field is equally distributed across academic labs, industrial research groups, and open-source communities. Stanford, Princeton, CMU, AI2, UCLA, UBC, Vector Institute, Peking University, Zhejiang University, Shanghai Jiao Tong University, Google Research, Google DeepMind, Anthropic, Microsoft, OpenAI, xAI, Alibaba, ModelScope, Sakana-style open-ended evolution groups, and many independent developers all contribute different pieces. Academic labs often contribute conceptual clarity and benchmarked methods. Industrial labs contribute scale, infrastructure, and high-value evaluation domains. Open-source projects contribute rapid experimentation, usability pressure, and failure reports from real developers.

This map matters because future progress will require cross-community synthesis. Researchers who study RL reward design need to talk to software engineers who understand CI and rollback. Evolutionary-computation researchers need to talk to security engineers who understand supply-chain risk. Memory researchers need to talk to privacy and governance experts. Product teams need to talk to benchmark designers so that production reliability is not reduced to static leaderboard scores. The field's researcher graph is already interdisciplinary; its evaluation and publication norms should become interdisciplinary as well.

\subsection{Open Problems That Define the Next Stage}

The most urgent open problem is evaluator reliability. Every loop depends on feedback, and every feedback signal can be gamed. LLM-as-a-judge is flexible but may be biased, unstable, or persuadable. Unit tests are precise but incomplete. Benchmarks are comparable but can be overfit. Human review is high quality but expensive and inconsistent. Future systems need evaluator ensembles, hidden tests, adversarial validation, evaluator versioning, and explicit reporting of evaluator limitations. A self-evolving agent should not be allowed to modify or select its final gate evaluator.

The second open problem is cost. Open-ended search can consume enormous compute, tokens, tool calls, human review time, and storage. A practical evolutionary agent must treat cost as part of the fitness function, not as an afterthought. It should report marginal gain per token, rollback rate, failed-candidate ratio, wall-clock time, and maintenance cost. Many production teams do not need the most open-ended system; they need the smallest loop that reliably improves a bounded workflow.

The third open problem is memory governance. Long-term memory is essential for cumulative learning, yet unmanaged memory becomes a liability. Systems need provenance, summarization, conflict resolution, time decay, privacy filtering, poisoning detection, and task-sensitive retrieval. Memory should not be a flat transcript dump. It should be an auditable knowledge base with versioned claims and links back to evidence.

The fourth open problem is safe self-modification. Code and architecture evolution require permissions, sandboxes, static checks, dynamic tests, canary deployments, and rollback. The agent should generate candidate changes, but the environment should decide which changes can execute and under what constraints. The stronger the mutation operator, the stronger the isolation boundary must be.

The fifth open problem is transfer and negative transfer. The field must distinguish benchmark-specific improvement from general adaptation. Every claimed evolution result should ask: what transfers, where does it fail, which old capabilities regress, and how is capability erosion detected? Regression suites, retained ancestor archives, task-slice reporting, and cross-domain evaluation are essential.

The sixth open problem is social accountability. Self-evolving agents will increasingly operate inside organizations. That means they need ownership, approval paths, audit logs, incident procedures, and human override. Agent evolution is not just an ML training problem; it is a software-change-management and governance problem.

\subsection{A Visionary but Grounded Outlook}

The most credible future for agent self-evolution is not a sudden jump to unconstrained artificial general intelligence. It is a staged expansion of governed open-endedness. In the near term, the strongest deployments will appear in domains with clear feedback: code maintenance, test generation, data analysis, document repair, workflow automation, web tasks with verifiable outcomes, scientific code optimization, and operations support. These domains allow agents to propose changes while independent infrastructure tests those changes. They also have enough economic value to justify the cost of evaluation.

In the medium term, agent evolution will likely become an infrastructure layer rather than a feature of individual agents. Teams will maintain evaluator libraries, lineage stores, memory banks, sandbox farms, cost dashboards, and policy gates. Agents will submit candidate improvements to this infrastructure. Low-risk changes will be automatically evaluated and merged; medium-risk changes will be staged; high-risk changes will require human approval. This resembles a future CI/CD system for intelligent behavior: not just continuous integration of code, but continuous integration of prompts, skills, tools, memories, policies, and workflows.

In the longer term, if transfer, evaluation, and governance mature together, agent self-evolution could become one of the engineering paths toward more general adaptive systems. The key phrase is ``together.'' Open-ended exploration without governance becomes unsafe. Governance without exploration becomes static. Evaluation without memory produces repeated rediscovery. Memory without forgetting produces drift. Architecture search without cost control produces beautiful but unusable complexity. The future discipline must coordinate all five loops under explicit safety, cost, and accountability constraints.

The introduction framed agent evolution as a bridge between machine learning, evolutionary computation, program synthesis, reinforcement learning, meta-learning, cognitive architecture, and software engineering. The conclusion strengthens that claim. The field's central insight is that improvement is no longer confined to model weights. It can occur in the surrounding system: in the prompt that frames a task, the evaluator that detects failure, the tool that expands action space, the memory that preserves experience, the workflow that coordinates roles, and the code that implements the agent itself. That expansion is why agent self-evolution is so promising and so risky.

A responsible research agenda should therefore pursue three commitments. First, make every improvement claim falsifiable: specify the mutable object, feedback source, selection rule, evaluator independence, cost, and failure cases. Second, make every persistent change reversible: preserve lineage, ancestor states, audit trails, and rollback paths. Third, make every open-ended system socially embedded: define who owns it, who can stop it, who reviews high-risk mutations, and whose values shape the evaluator. If the community follows these commitments, agent self-evolution can move beyond isolated demonstrations toward systems that genuinely learn from interaction while remaining inspectable, corrigible, and useful.

The final judgment of this survey is therefore optimistic but conditional. The evidence from papers, projects, repositories, user pain points, and researcher networks shows that agent self-evolution is becoming a real field. It already has recurring loops, representative systems, emerging infrastructure, and a shared set of unsolved problems. But its success will not be measured by whether an agent can claim to improve itself. It will be measured by whether the improvement survives independent evaluation, transfers beyond the training niche, respects safety boundaries, reduces human toil, and accumulates knowledge that future agents and human teams can trust. The path forward is not evolution instead of engineering; it is evolution made into engineering.
